\documentclass[11pt]{article}
\usepackage[margin=0.95in]{geometry}
\usepackage{amsmath,amssymb}
\usepackage{booktabs}
\usepackage{array,float}
\newcolumntype{L}[1]{>{\raggedright\arraybackslash}p{#1}}
\usepackage{enumitem}
\usepackage{microtype}
\usepackage[hidelinks]{hyperref}

\setlist{nosep,leftmargin=*}
\setlength{\parskip}{0pt}
\setlength{\parindent}{0pt}
\makeatletter
\renewcommand\paragraph{\@startsection{paragraph}{4}{\z@}%
  {2.5pt\@plus1pt\@minus.2pt}{-0.5em}{\normalfont\normalsize\bfseries}}
\renewcommand\section{\@startsection{section}{1}{\z@}%
  {7pt\@plus2pt\@minus2pt}{2pt}{\normalfont\large\bfseries}}
\makeatother

\title{\vspace{-1.5em}\textbf{Literature Review: Conditional Generation under Weak Conditioning}\\[4pt]
\large Independent Study --- Deep Learning}
\author{Zhizhe Zhang (zz5070) \quad Dixuan Yang (dy2397)\\
NYU Shanghai \quad$\cdot$\quad Advisor: Prof.\ Shengjie Wang}
\date{September 2026}

\begin{document}
\maketitle
\vspace{-3em}

\section{Scope}

Our project asks how conditional generative models behave when the conditioning signal is weak, using EEG-conditioned hand-trajectory generation as the testbed. We review the task we adopt and intend to audit (\S2), two further reference points on the same data (\S3), conditional generation from weak neural conditions in other modalities (\S4), and the generative machinery plus the condition-fidelity metrics that some output domains have and the EEG literature does not (\S5).

\section{Continuous hand kinematics from EEG}

Nine studies attack essentially one task --- regress continuous 3D hand kinematics from low-frequency scalp EEG, scored by Pearson correlation (PCC) --- and six use the dataset we plan to use. Table~\ref{tab:lit} gives each one's task, method, result, and whether it computed a chance level.

\begin{table}[H]
\centering\footnotesize
\setlength{\tabcolsep}{4pt}
\begin{tabular}{@{}L{2.15cm}L{6.2cm}L{3.5cm}L{3.0cm}@{}}
\toprule
Study & Task and method & Reported PCC $x/y/z$ & Chance control \\
\midrule
Bradberry 2010 & 3D center-out reaching, 5 subjects; linear decoder on ${<}1$\,Hz EEG amplitudes, multi-lag, sLORETA sources & ``comparable to intracranial''; falls as movement variability rises & none \\
Korik 2018 & \emph{imagined} 3D trajectories; mLR on mu/beta/low-gamma band power & band power beats filtered potentials & none \\
Jain \& Kumar 2022/2023 & WAY-EEG-GAL, pre-movement delta EEG (lags 150--350\,ms); CNN--LSTM; LOSO added in 2023 & $0.79/0.80/0.63$; \textbf{LOSO $0.76/0.78/0.48$} & none \\
ESI-GAL 2024 & WAY-EEG-GAL; frontoparietal source imaging $+$ rEEGNet & $0.790/0.795/0.637$ sensor; $0.769/0.777/0.647$ source & none \\
M3T-Attention 2026 & WAY-EEG-GAL; multi-scale temporal attention transformer, 0.5--12\,Hz $\to$ pos./vel./accel. & $0.8816/0.8841/0.8711$ & none; splits by session within one subject \\
SAND 2026 & WAY-EEG-GAL; spectral $+$ temporal-attention dual branch, five-fold CV & $0.9595/0.9534/0.9293$ & none \\
\midrule
Antelis 2013 & linear regression from low-frequency EEG vs.\ randomly re-paired data & \emph{not above chance} & re-pairing \\
Kobler 2018 & visuomotor/oculomotor tracking; shuffle null plus a hand-at-rest replay control & $0.2$--$0.4$ (chance $0.10$--$0.13$; hand at rest $0.33$--$0.36$) & shuffle $+$ behavioural \\
Hosseini 2022 & hand movement from phase-based connectivity (PLV, MSC) & $0.43\pm0.03$ & two surrogates \\
\bottomrule
\end{tabular}
\caption{EEG-to-kinematics studies. The three that compute a chance level sit below the rule.}
\label{tab:lit}
\end{table}

\vspace{-0.7em}
\paragraph{Relevance of each.} \emph{Bradberry} is the founding claim we stress-test, and its finding that accuracy falls as movement variability rises is the first warning sign, being equally what a model fitting the task-average trajectory produces. \emph{Korik} extends it to imagery but inherits the metric and the missing null. \emph{Jain \& Kumar} is our closest baseline in task definition and contributes the most diagnostic number here: LOSO sits within $0.03$--$0.15$ of within-subject performance, yet EEG is strongly subject-idiosyncratic, so a decoder transferring this well to unseen subjects is more plausibly reproducing the task-average trajectory than reading subject-specific cortical activity. \emph{ESI-GAL} shows source imaging from 32 electrodes gives no gain over raw sensors. \emph{M3T-Attention} is our reproduction target, its code being public. \emph{SAND} ends a four-year climb on one dataset ($\approx0.5\to0.79\to0.88\to0.96$) across 3{,}936 nearly identical trials; $0.96$ between scalp EEG and hand position would exceed much of the intracranial literature.

\paragraph{The pattern.} \emph{Antelis} explains why: PCC is scale-invariant and systematically elevated for slow, sinusoid-like signals, which low-pass EEG and hand velocity both are --- hence displacement error in millimetres as our primary metric. \emph{Kobler} quantifies how much correlation arises with no movement at all, and shows cross-subject averaging raises $r$ from $0.40$ to $0.83$, hence our per-subject distributions. Together: \textbf{studies that compute a chance level report ${\approx}0.4$; studies that do not report ${\approx}0.9$.}

\section{Two further reference points}

\paragraph{Combrisson \& Jerbi (2015).} Analyses the null distribution of decoding accuracy at finite sample sizes, shows theoretical chance levels are anticonservative, and supplies permutation thresholds. \emph{Relevance:} justifies defining Gate 1 against an empirical null rather than zero.

\paragraph{Barachant \& Bonnet (2015), Kaggle Grasp-and-Lift.} Detects six discrete grasp-and-lift events sample-wise from the same EEG and same 12 subjects using Riemannian covariance features ensembled with CNNs and RNNs: 0.981 AUC, first of 379 teams. \emph{Relevance:} on identical data the discrete problem is solved to near ceiling under adversarial evaluation by hundreds of teams, while regression never faced one --- the asymmetry underwriting our fallback of coarser output if Gate 1 fails.

\section{Conditional generation from weak neural conditions}

\paragraph{Takagi \& Nishimoto (2023); Shirakawa et al.~(2025).} The first reconstructs viewed images from fMRI via a pretrained diffusion model's latent and text-embedding spaces, scored by per-image quality; the second tests whether such reconstructions carry \emph{sample-specific} information, using data built to prevent train/test category overlap, and finds performance collapses out of distribution --- the results being largely classification into trained categories plus hallucination by the prior. \emph{Relevance:} the failure we predict for EEG-to-motion, already shown in a better-resourced modality and invisible to per-sample quality metrics.

\paragraph{D\'efossez et al.~(2023).} Decodes perceived speech from MEG/EEG (175 subjects, $>$160\,h) by contrastive alignment with wav2vec 2.0 embeddings; from 3\,s of MEG the correct segment is identified among $>$1{,}000 candidates with up to 41\% accuracy. \emph{Relevance:} the key precedent --- under weak conditioning this field evaluates by \emph{identification among candidates}, which our cross-conditioning matrix generalizes to a domain with no target encoder.

\paragraph{NeuroSonic (Gao et al., 2026).} Reconstructs speech from scalp EEG by conditional flow matching --- deterministic probability-flow velocity field, $x$-prediction with a velocity-space loss, EEG and audio in a shared token space --- gaining up to 26.3\% perceptual quality on CineBrain and EAV. \emph{Relevance:} our closest methodological neighbour, but its metrics (FAD, LSD, spectral convergence, DNSMOS) never depend on \emph{which} EEG trial was the condition, and with no shuffle control, template baseline or guidance analysis, a model ignoring EEG and generating plausible speech would score well on every reported number.

\paragraph{JET (Wang et al., 2026).} Generates EEG for augmentation by conditional flow matching with structure-preserving constraints, cutting TS-FID by over 40\%; ablations show reconstruction loss alone is worst, constraint weight has a concave optimum, and a zero-noise base is catastrophic. \emph{Relevance:} supplies our loss design, our Gaussian base, and one diagnostic (silhouette of samples clustered by condition) --- the only condition-fidelity measure in this literature.

\paragraph{FDA (Wang et al., 2025).} Few-trial adaptation of neural-to-behaviour decoders on invasive primate recordings using flow matching with stable latent dynamics. \emph{Relevance:} flow matching works for neural-to-behaviour mapping in our low-data regime, though with a far stronger condition.

\section{Generative machinery and condition-fidelity metrics}

\paragraph{Flow matching and guidance.} Lipman et al.~(2023) give the flow-matching objective, Ho \& Salimans (2022) classifier-free guidance, Peebles \& Xie (2023) the AdaLN Transformer, Li \& He (2025) the $x$-prediction argument --- together, our generator. Kynk\"a\"anniemi et al.~(2024) show CFG is harmful at high noise levels and useful only in a middle interval (ImageNet-512 FID $1.81\!\to\!1.40$); Karras et al.~(2024) guide with a weakened unconditional model instead. \emph{Relevance:} the guidance scale is not transferable even under strong conditioning, and with a noisy condition the difference CFG amplifies is largely noise --- making our CFG sweep a first-class experiment.

\paragraph{Motion generation, R-precision, best-of-$K$.} Tevet et al.~(2023, MDM) generate motion from text with a diffusion transformer predicting the sample rather than the noise, so geometric losses on positions and velocities apply --- the parameterization we adopt for the same reason --- while Guo et al.~(2022) contribute \textbf{R-precision}, how often a generated motion retrieves its own text among 32 candidates. \emph{Relevance:} motion generation already has a condition-fidelity metric, and it is a retrieval metric --- the same construct as D\'efossez's identification accuracy and our cross-conditioning matrix. The counter-example is Gupta et al.~(2018), whose min-ADE over $K\!=\!20$ samples rewards covering plausible futures and cannot separate a model using its condition from one sampling broadly --- the trap our protocol forbids.

\paragraph{Representation and condition collapse.} Ijspeert et al.~(2013) formalize dynamic movement primitives, whose arc-length variant decouples path geometry from execution timing --- our structured output, claimed not as novel but for its role under weak conditioning, where lower output dimensionality means less information the condition must supply. \emph{Relevance:} the failure we study also has a name outside BCI --- posterior collapse in conditional VAEs (Bowman et al., 2016).

\section{Datasets and tasks planned}

\paragraph{WAY-EEG-GAL (Luciw, Jarocka \& Edin, 2014).} Twelve right-handed participants (8 female, 19--35). Each trial: LED cue, reach, grasp with thumb and index finger, lift, hold ${\approx}2$\,s, replace, release, return to rest, with object weight (165/330/660\,g) and surface friction (sandpaper/suede/silk) changing unpredictably between trials. Ten series per participant (6 weight series of 34 lifts, 2 friction of 34, 2 mixed of 28) give 328 trials each, 3{,}936 total, in 12 archives of 10.34\,GB. Recorded: 32-channel actiCAP/BrainAmp EEG at 500\,Hz, standard 10--20 montage, no artifact rejection; 5-channel EMG at 4\,kHz; a Polhemus FASTRAK with four 6-DOF sensors on object, index finger, thumb and \textbf{wrist}; two ATI Nano17 force/torque sensors --- non-EEG streams nominally 500\,Hz, synchronized within 2\,ms. Sixteen event times and 43 per-trial measures (\texttt{tHandStart}, \texttt{tPeakVelHandReach}, \texttt{Dur\_Reach}) are released, which we use for onset alignment rather than our own detector.

Three properties shape our plan. \emph{(i) Usable trials are 3{,}528, not 3{,}936}: the non-EEG streams of one weight series per participant were withheld for a later competition (which became the Kaggle challenge), leaving 294 multimodal trials per subject, so holding out 3--4 subjects leaves 2{,}352--2{,}646 training pairs. \emph{(ii) Kinematics are nominally 500\,Hz but the hardware cannot deliver it}: FASTRAK updates at 120\,Hz divided by active sensors, i.e.\ ${\approx}30$\,Hz with four, with no interpolation documented, so the true resolution of our target bounds every timing metric we report. \emph{(iii) Coverage is sparse with no EOG channel}: contralateral sensorimotor cortex is covered by roughly C3, Cz, FC1, FC5, CP1 and CP5, while Fp1/Fp2 are present and unrejected, making ocular contamination a live hypothesis; TP9/TP10 are data channels, so mastoid re-referencing remains available.

\paragraph{Task and baselines.} Input: ${\approx}1$\,s of EEG preceding movement onset, and nothing else. Output: the full 3D wrist displacement trajectory of the subsequent reach (1--2\,s), in one pass, relative to onset position. Baselines under identical splits: a phase-conditioned template ignoring its input; a shuffled-pairing model; an EMG-conditioned positive control; a lightweight EEG regressor; a reproduction of M3T-Attention. A second testbed degrades a clean motion-generation condition along three axes --- additive noise, an information bottleneck, per-``subject'' shift --- so condition strength becomes an independent variable and EEG one point on a curve; it is not data-limited, and is the fallback if Gate 1 fails.

\section{Hypotheses to test over the next two weeks}

Gate 1 asks whether EEG carries usable trajectory information above an empirically estimated null, on the exact dataset and split we intend to use.

\begin{enumerate}
\item \textbf{Template.} A phase-conditioned mean-trajectory model receiving no EEG input reaches per-axis PCC within 0.1 of published values ($\gtrsim0.8$ on $x,y$). \emph{If confirmed,} PCC on this dataset is uninformative about condition usage and the published ranking of decoders is not evidence of decoding.
\item \textbf{Shuffle.} A decoder trained on randomly re-paired EEG and trajectories is not separable from a correctly paired one by PCC, but is separable by mean displacement error; 50--100 retrained shuffled models give the empirical null under both, fixing which metric we use thereafter.
\item \textbf{Metric reversal.} The rank order of the five baselines changes when PCC is replaced by mean 3D displacement error in millimetres.
\item \textbf{Cross-subject anomaly.} For the reproduced M3T, LOSO performance falls within 0.05 PCC of within-subject performance --- which would mean it uses no subject-specific neural information, and that published LOSO results are template-fitting.
\item \textbf{Ocular.} Removing Fp1, Fp2, F7 and F8 does not push the EEG model below the template level; if performance instead collapses to chance, the signal was ocular rather than cortical.
\end{enumerate}

Two measurements precede these: the effective sampling rate of the FASTRAK kinematics, and the usable trial count with the distribution of \texttt{Dur\_Reach}. \emph{Gate 1 pass criterion, fixed in advance:} the EEG model's mean displacement error on held-out subjects falls below the 5th percentile of the shuffle distribution. We expect hypotheses 1, 2 and 4 to be confirmed and the criterion met only narrowly --- which would establish the project's premise: real but small EEG information, reported by the standard metrics as large.

\begin{thebibliography}{99}\small
\bibitem{antelis2013} J.~M. Antelis, L.~Montesano, A.~Ramos-Murguialday, N.~Birbaumer, J.~Minguez. On the usage of linear regression models to reconstruct limb kinematics from low frequency EEG signals. \emph{PLoS ONE}, 8(4):e61976, 2013.
\bibitem{barachant2015} A.~Barachant, R.~Bonnet. Winning solution of the Grasp-and-Lift EEG Detection challenge (Kaggle; 379 teams, 0.981 AUC), 2015. \url{https://github.com/alexandrebarachant/Grasp-and-lift-EEG-challenge}.
\bibitem{bowman2016} S.~R. Bowman, L.~Vilnis, O.~Vinyals, A.~M. Dai, R.~J\'ozefowicz, S.~Bengio. Generating sentences from a continuous space. \emph{CoNLL}, 2016.
\bibitem{bradberry2010} T.~J. Bradberry, R.~J. Gentili, J.~L. Contreras-Vidal. Reconstructing three-dimensional hand movements from noninvasive electroencephalographic signals. \emph{J.\ Neuroscience}, 30(9):3432--3437, 2010.
\bibitem{combrisson2015} E.~Combrisson, K.~Jerbi. Exceeding chance level by chance: the caveat of theoretical chance levels in brain signal classification and statistical assessment of decoding accuracy. \emph{J.\ Neuroscience Methods}, 250:126--136, 2015.
\bibitem{defossez2023} A.~D\'efossez, C.~Caucheteux, J.~Rapin, O.~Kabeli, J.-R. King. Decoding speech perception from non-invasive brain recordings. \emph{Nature Machine Intelligence}, 5:1097--1107, 2023.
\bibitem{sand2026} R.~Fu, Y.~Fang, F.~Xu, C.~Hua, C.~Hua. SAND: spectral-attention neural decoding of hand kinematics from low-frequency EEG dynamics. \emph{IEEE Trans.\ Biomedical Engineering}, 2026.
\bibitem{neurosonic} W.~Gao, Y.~Wang, Y.~Ma, C.~Yang, W.~Li, C.~You. NeuroSonic: conditional flow matching for EEG-to-speech reconstruction. \emph{MICCAI}, 2026. arXiv:2606.24087.
\bibitem{guo2022} C.~Guo, S.~Zou, X.~Zuo, S.~Wang, W.~Ji, X.~Li, L.~Cheng. Generating diverse and natural 3D human motions from text. \emph{CVPR}, 2022.
\bibitem{sgan2018} A.~Gupta, J.~Johnson, L.~Fei-Fei, S.~Savarese, A.~Alahi. Social GAN: socially acceptable trajectories with generative adversarial networks. \emph{CVPR}, 2018.
\bibitem{cfg2022} J.~Ho, T.~Salimans. Classifier-free diffusion guidance. arXiv:2207.12598, 2022.
\bibitem{hosseini2022} S.~M. Hosseini, V.~Shalchyan. Continuous decoding of hand movement from EEG signals using phase-based connectivity features. \emph{Frontiers in Human Neuroscience}, 16:901285, 2022.
\bibitem{dmp2013} A.~J. Ijspeert, J.~Nakanishi, H.~Hoffmann, P.~Pastor, S.~Schaal. Dynamical movement primitives: learning attractor models for motor behaviors. \emph{Neural Computation}, 25(2):328--373, 2013. See also: Arc-length dynamic movement primitives, \emph{Robotics and Autonomous Systems}, 2017.
\bibitem{premovnet} A.~Jain, L.~Kumar. PreMovNet: pre-movement EEG-based hand kinematics estimation for grasp-and-lift task. \emph{IEEE Sensors Letters}, 2022. arXiv:2205.01050.
\bibitem{jain2023} A.~Jain, L.~Kumar. Subject-independent trajectory prediction using pre-movement EEG during grasp and lift task. \emph{Biomedical Signal Processing and Control}, 86:105160, 2023. arXiv:2209.01932.
\bibitem{esigal} A.~Jain, L.~Kumar. ESI-GAL: EEG source imaging-based kinematics parameter estimation for grasp and lift task. \emph{Computers in Biology and Medicine}, 2024. arXiv:2406.11500.
\bibitem{autoguidance} T.~Karras, M.~Aittala, T.~Kynk\"a\"anniemi, J.~Lehtinen, T.~Aila, S.~Laine. Guiding a diffusion model with a bad version of itself. \emph{NeurIPS}, 2024.
\bibitem{kobler2018} R.~J. Kobler, A.~I. Sburlea, G.~R. M\"uller-Putz. Tuning characteristics of low-frequency EEG to positions and velocities in visuomotor and oculomotor tracking tasks. \emph{Scientific Reports}, 8:17713, 2018.
\bibitem{korik2018} A.~Korik, R.~Sosnik, N.~Siddique, D.~Coyle. Decoding imagined 3D hand movement trajectories from EEG: evidence to support the use of mu, beta, and low gamma oscillations. \emph{Frontiers in Neuroscience}, 12:130, 2018.
\bibitem{guidanceinterval} T.~Kynk\"a\"anniemi, M.~Aittala, T.~Karras, S.~Laine, T.~Aila, J.~Lehtinen. Applying guidance in a limited interval improves sample and distribution quality in diffusion models. \emph{NeurIPS}, 2024.
\bibitem{jit} T.~Li, K.~He. Back to basics: let denoising generative models denoise. arXiv:2511.13720, 2025.
\bibitem{fm} Y.~Lipman, R.~T.~Q. Chen, H.~Ben-Hamu, M.~Nickel, M.~Le. Flow matching for generative modeling. \emph{ICLR}, 2023.
\bibitem{wayeeggal} M.~D. Luciw, E.~Jarocka, B.~B. Edin. Multi-channel EEG recordings during 3,936 grasp and lift trials with varying weight and friction. \emph{Scientific Data}, 1:140047, 2014.
\bibitem{dit} W.~Peebles, S.~Xie. Scalable diffusion models with transformers. \emph{ICCV}, 2023.
\bibitem{shirakawa2025} K.~Shirakawa, Y.~Nagano, et al. Spurious reconstruction from brain activity. \emph{Neural Networks}, 2025. arXiv:2405.10078.
\bibitem{takagi2023} Y.~Takagi, S.~Nishimoto. High-resolution image reconstruction with latent diffusion models from human brain activity. \emph{CVPR}, 2023.
\bibitem{mdm} G.~Tevet, S.~Raab, B.~Gordon, Y.~Shafir, D.~Cohen-Or, A.~H. Bermano. Human motion diffusion model. \emph{ICLR}, 2023.
\bibitem{fda} P.~Wang et al. Flow matching for few-trial neural adaptation with stable latent dynamics. \emph{ICML}, 2025.
\bibitem{jet} Y.~Wang, Y.~Ma, W.~Li, C.~You. Let EEG models learn EEG. \emph{ICML}, 2026. arXiv:2605.21280.
\bibitem{m3t} L.~Zhu, P.~Jiang, A.~Huang, J.~Zhang, P.~Yuan. M3T-attention: a multi-level multi-scale temporal attention transformer for EEG hand movement trajectory decoding. \emph{Cognitive Neurodynamics}, 20(1):33, 2026. Code: \url{https://github.com/jjspp/M3T_Attention}.
\end{thebibliography}

\end{document}
